\documentclass[11pt,a4paper]{article}
\usepackage[utf8]{inputenc}
\usepackage[T1]{fontenc}
\usepackage[spanish]{babel}
\usepackage[a4paper,margin=2.4cm]{geometry}
\usepackage{booktabs}
\usepackage{array}
\usepackage{xcolor}
\usepackage{listings}
\usepackage{hyperref}

\lstset{basicstyle=\ttfamily\small,breaklines=true,frame=single,columns=fullflexible,showstringspaces=false}
\hypersetup{colorlinks=true,linkcolor=blue,urlcolor=blue}

\title{\textbf{Modelo de datos NoSQL / especializado}\\[0.3em]
\Large TimescaleDB + PostGIS para predicción de demanda}
\author{García, Barreto, Chulluncuy, Hilario}
\date{Base de Datos (MCD8008) -- Parte 2}

\begin{document}
\maketitle

\section{Principios de diseño}

El modelo especializado \textbf{no reemplaza} al modelo relacional. Se trata de una capa
analítica que reutiliza las dimensiones (\texttt{core.zona\_taxi}, \texttt{core.plataforma})
y proyecta la serie de demanda en una estructura optimizada para consultas temporales y
espaciales.

\begin{itemize}
  \item \textbf{Polyglot Persistence:} PostgreSQL conserva el núcleo transaccional; la capa
        TimescaleDB/PostGIS se deriva de él y se reconstruye desde el agregado
        \texttt{analytics.demanda\_hora}.
  \item \textbf{Desnormalización deliberada y reconstruible:} el hypertable duplica la demanda
        horaria para acelerar el acceso, igual que ya hacía \texttt{demanda\_hora} en la Parte 1.
  \item \textbf{Separación temporal--espacial:} las coordenadas viven en la entidad maestra
        \texttt{zona\_taxi} (normalizada) y se unen en tiempo de consulta; así el hypertable
        permanece compresible (las columnas geográficas no son compresibles de forma nativa).
\end{itemize}

\section{Capa espacial (PostGIS)}

\subsection{Atributo geo en la entidad zona\_taxi}
\begin{lstlisting}
ALTER TABLE core.zona_taxi
    ADD COLUMN IF NOT EXISTS geo GEOGRAPHY(POINT, 4326);

UPDATE core.zona_taxi
SET geo = ST_SetSRID(ST_MakePoint(longitud, latitud), 4326)::geography
WHERE longitud IS NOT NULL AND latitud IS NOT NULL;

CREATE INDEX IF NOT EXISTS idx_zona_taxi_geo
    ON core.zona_taxi USING GIST (geo);
\end{lstlisting}

Se usa el tipo \texttt{GEOGRAPHY(POINT,4326)} (WGS84) para que las funciones de distancia
(\texttt{ST\_Distance}, \texttt{ST\_DWithin}) devuelvan valores en metros sobre el elipsoide,
adecuados para el dominio de NYC.

\section{Hypertable de demanda (TimescaleDB)}

\subsection{Estructura}
\begin{lstlisting}
CREATE TABLE analytics.demanda_serie (
    hora            TIMESTAMPTZ NOT NULL,
    zona_origen_id  SMALLINT    NOT NULL,
    plataforma_codigo VARCHAR(10) NOT NULL,
    viajes          BIGINT      NOT NULL
);

SELECT create_hypertable('analytics.demanda_serie', 'hora',
                         if_not_exists => TRUE);
\end{lstlisting}

La columna \texttt{hora} es la dimensión de particionamiento: TimescaleDB la divide en
\emph{chunks} de tamaño automático, de modo que una consulta por rango temporal solo accede
a los chunks relevantes.

\subsection{Población}
\begin{lstlisting}
INSERT INTO analytics.demanda_serie
    (hora, zona_origen_id, plataforma_codigo, viajes)
SELECT hora::timestamptz, zona_origen_id, plataforma_codigo, viajes
FROM analytics.demanda_hora;
\end{lstlisting}

\section{Continuous aggregate}

\begin{lstlisting}
CREATE MATERIALIZED VIEW analytics.demanda_diaria
WITH (timescaledb.continuous) AS
SELECT time_bucket('1 day', hora) AS dia,
       zona_origen_id,
       plataforma_codigo,
       SUM(viajes) AS viajes,
       COUNT(*)     AS n_horas
FROM analytics.demanda_serie
GROUP BY dia, zona_origen_id, plataforma_codigo
WITH NO DATA;

CALL refresh_continuous_aggregate('analytics.demanda_diaria', NULL, NULL);
\end{lstlisting}

El agregado diario se mantiene de forma incremental y automática, evitando recalcular
agregaciones completas en cada consulta.

\section{Política de compresión}

\begin{lstlisting}
ALTER TABLE analytics.demanda_serie SET (
    timescaledb.compress,
    timescaledb.compress_segmentby = 'zona_origen_id, plataforma_codigo',
    timescaledb.compress_orderby   = 'hora'
);

SELECT add_compression_policy('analytics.demanda_serie', INTERVAL '7 days');
\end{lstlisting}

Los segmentos se comprimen por \texttt{zona\_origen\_id} y \texttt{plataforma\_codigo}, y se
ordenan por \texttt{hora}. Resultado observado: $162$\,MB (relacional) $\to$ $14$\,MB
(comprimido), factor $\approx 11{,}6\times$.

\section{Esquema resultante}

\begin{center}
\begin{tabular}{@{}ll@{}}
\toprule
\textbf{Objeto} & \textbf{Tecnología / tipo} \\
\midrule
\texttt{core.zona\_taxi.geo} & PostGIS \texttt{GEOGRAPHY(POINT,4326)} + índice GiST \\
\texttt{analytics.demanda\_serie} & Hypertable (dimensión \texttt{hora}) \\
\texttt{analytics.demanda\_diaria} & Continuous aggregate diario \\
Compresión & Segmento por (zona, plataforma), orden por hora \\
\bottomrule
\end{tabular}
\end{center}

\end{document}
